\documentclass[11pt]{article}
\usepackage[margin=1in]{geometry}
\usepackage{amsmath,amssymb,amsthm,mathtools}
\usepackage[hidelinks]{hyperref}
\newtheorem{theorem}{Theorem}
\newtheorem{lemma}[theorem]{Lemma}
\newtheorem{corollary}[theorem]{Corollary}
\newcommand{\C}{\mathbb C}
\newcommand{\Q}{\mathbb Q}
\newcommand{\Qb}{\overline{\mathbb Q}}
\newcommand{\Lfac}{\mathcal L}
\DeclareMathOperator{\Res}{Res}
\DeclareMathOperator{\Ann}{Ann}
\title{Simplex pencil phases and affine radial limits in every dimension}
\author{Internal proof and exact implementations; no priority claim}
\date{16 November 2026}
\begin{document}
\maketitle
\noindent\textbf{Scope and credit.}
This note proves a whole class in every dimension: polynomial phases on a
simplex that depend on one complex affine linear form, and homogeneous
factorial phases in the radial sum and one other linear form. Repeated
coefficients and arbitrary polynomial markers are retained. A separate
residue argument gives nonvanishing of affine exponential simplex periods
at algebraic coefficients, with an application to arbitrary lower layers
of flat factorial polynomials.

The divided-difference identities are classical Genocchi--Hermite formulas;
see \href{https://arxiv.org/abs/math/0502036}{de Boor, Divided Differences
(2005), especially the Genocchi--Hermite section}. The path/Cauchy argument
is the polynomial moment mechanism used in
\href{https://arxiv.org/abs/0710.4085v2}{Pakovich--Muzychuk} and the
\href{sources/long-eit-pinned.tex}{pinned EIT I tree-density proof}.
It is written here for finite sums with different polynomial amplitudes.
The flat radial limit is the all-dimensional theorem in
\href{sources/long-eit-vi-current.tex}{Long--Le Blanc, EIT VI, Theorem 2.1}.
No theorem from EIT V or any arithmetic EIT lifting theorem is needed for
the pencil assertion. This is an internal written result, with no claim
of novelty, outside refereeing, or formal verification.

\section{A polynomial-chain functional cannot have a nonzero finite tail}
Let
\[
 \Lambda(F)=\sum_{j=1}^s\int_{a_j}^{b_j} r_j(z)F(z)\,dz,
 \qquad r_j\in\C[z],\quad a_j,b_j\in\C.                    \tag{1}
\]
The paths may be chosen arbitrarily because every integrand is entire
when $F$ is entire. Complex linear coefficients can be absorbed in $r_j$.

\begin{lemma}\label{lem:chain}
If $p\in\C[z]$ is nonconstant and $0\ne H\in\C[t]$ satisfies
$\Lambda(H(p)p^k)=0$ for every $k\ge0$, then
$\Lambda(F(p))=0$ for every entire $F$. In particular, if
$\Lambda(p^m)=0$ eventually, then $\Lambda(p^m)=0$ for every $m\ge0$.
\end{lemma}
\begin{proof}
Choose an embedded finite tree $T$ containing all finite critical values
of $p$ and the values of every endpoint in (1). The inverse graph
$p^{-1}(T)$ is connected: after adding a regular central vertex and arms,
local monodromies generate the transitive polynomial covering group.
Consequently each pair of endpoints can be connected on this graph.
Replace every path in (1) by such a graph path. Polynomial and entire
integrands keep the same integral, separately for each $r_j$.

Orient and subdivide the target edges. Over an open edge $e$ the combined
pushforward density has the form
\[
 \psi_e(t)=\sum_{j,i} n_{j,e,i}
              \frac{r_j(z_i(t))}{p'(z_i(t))}.
\]
It is analytic there and has only integrable endpoint singularities of
orders at worst $1/k-1>-1$. For $u\notin T$ form
\[
 C_H(u)=\sum_e\int_e\frac{H(t)\psi_e(t)}{u-t}\,dt.
\]
Its Laurent expansion at infinity has the assumed zero moments as its
coefficients, so it is zero near infinity. The complement of a finite
tree is connected, and analytic continuation gives zero on that complement.
Across an interior edge its jump is $\pm2\pi i H(t)\psi_e(t)$; all other
edge contributions extend holomorphically locally. Thus $\psi_e=0$
outside the finitely many zeros of $H$, and analyticity makes it identically
zero on every edge. Integrating $F(t)\psi_e(t)$ proves the first assertion.
For an eventually zero sequence take $H(t)=t^M$, for a sufficiently large
integer $M\ge1$.
\end{proof}

If $\Lambda(1)\ne0$, it follows that every nonconstant $p$ has infinitely
many nonzero moments. A nonzero constant $p$ has the same property.
The assumption that (1) has no point-evaluation terms matters:
$F\mapsto F(0)$ has mass one and annihilates every positive power of $z$.

\section{The simplex functional has a polynomial-chain representation}
Write $\Delta_{N-1}=\{t_i\ge0,\sum_i t_i=1\}$ with coordinate measure
of volume $1/(N-1)!$. Let $c_1,\ldots,c_N\in\C$ and
$\ell(t)=\sum_i c_i t_i$.

\begin{lemma}\label{lem:chain-representation}
If the $c_i$ are not all equal, the functional
$F\mapsto\int_{\Delta_{N-1}}F(\ell(t))dt$ has the form (1) with fixed
polynomial amplitudes of degrees at most $N-2$. This includes repeated nodes.
\end{lemma}
\begin{proof}
Put $k=N-1\ge1$ and take the entire $k$-fold primitive $G$ of $F$ whose
first $k$ jets at zero vanish. Genocchi--Hermite and the contour formula
for confluent divided differences give
\[
 \int_{\Delta_{N-1}}F(\ell(t))dt
   =[c_1,\ldots,c_N]G
   =\sum_{\lambda}\Res_{z=\lambda}\frac{G(z)}{Q(z)},
 \qquad Q(z)=\prod_i(z-c_i).                               \tag{2}
\]
These identities also follow termwise on monomials: the simplex integral
of $\ell^m$ is $m!h_m(c)/(m+N-1)!$, and the divided difference of
$m!z^{m+N-1}/(m+N-1)!$ is the same quantity. Uniform convergence on compact
sets then gives the entire-function identity.

If $\lambda$ has multiplicity $m_\lambda$ in $Q$, write
$1/Q=\sum_{\lambda,j} A_{\lambda,j}(z-\lambda)^{-j}$.
Its contribution to (2) is
$\sum_{j=1}^{m_\lambda}A_{\lambda,j}G^{(j-1)}(\lambda)/(j-1)!$.
There are at least two distinct nodes, so $m_\lambda\le N-1=k$.
Thus no term differentiates $G$ all the way to $F$. Instead,
\[
 G^{(j-1)}(\lambda)=
 \int_0^\lambda\frac{(\lambda-z)^{N-j-1}}{(N-j-1)!}F(z)dz.
\]
Substitution gives the explicit amplitudes
\[
 r_\lambda(z)=\sum_{j=1}^{m_\lambda}
 \frac{A_{\lambda,j}(\lambda-z)^{N-j-1}}
 {(j-1)!(N-j-1)!}.                                       \tag{3}
\]
All exponents are nonnegative. When all nodes coincide, (2) is instead
$F(c_1)/(N-1)!$, and must be treated as a point evaluation.
\end{proof}

\begin{theorem}[One affine coordinate on any simplex]\label{thm:simplex}
Let $\alpha_i$ be positive integers and $T\sim\operatorname{Dir}(\alpha)$.
For every complex affine linear form $\ell$ on the simplex and every
$p\in\C[z]$, if
\[
 \mathbb E[p(\ell(T))^m]=0\quad\hbox{eventually},
\]
then $p(\ell)$ is the zero polynomial function on the simplex.
If $p(\ell)$ is nonzero, then for every polynomial marker $w(T)$,
eventual vanishing of $\mathbb E[w(T)p(\ell(T))^m]$ implies its vanishing
for every $m\ge0$.
\end{theorem}
\begin{proof}
Absorb the affine constant into the coefficients of $\ell$ using
$\sum T_i=1$. For integer shapes, group $N=\sum\alpha_i$ uniform
simplex coordinates into blocks of sizes $\alpha_i$. Their sums have
the stated Dirichlet distribution, and $\ell$ has nodes $c_i$ repeated
$\alpha_i$ times. The normalized expectation is $(N-1)!$ times (2).
If $\ell$ is nonconstant its nodes are not all equal, so the preceding
lemma and Lemma~\ref{lem:chain} apply. The mass is one, which rules out
an eventually zero pure sequence for a nonzero phase.

For a polynomial marker, expand it into monomials $T^b$. The exact identity
\[
 \mathbb E_{\alpha}[T^b F(T)]
 =\frac{\prod_i(\alpha_i)_{b_i}}{(N)_{|b|}}
   \mathbb E_{\alpha+b}[F(T)]                             \tag{4}
\]
gives a finite sum of chain representations with the same nodes and
possibly larger integer multiplicities. Lemma~\ref{lem:chain} applies to
their sum even if its mass is zero. If $\ell$ or $p$ is constant, a
nonzero constant phase has moments equal to a nonzero constant power
times the marker's mean, so the same conclusion holds directly. The
zero phase is excluded in the marked assertion.
\end{proof}

In particular, every nonzero polynomial in a single complex affine
coordinate has a nonzero pure moment on a triangle, tetrahedron, or any
higher-dimensional simplex. Its image need not lie in a complex line:
three noncollinear complex node values already give image interior.
Polynomial marker periods cannot have a finite nonzero prefix followed
by an infinite zero tail for such a nonzero phase.

The same statement holds for uniform integration on a compact polytope,
and for polynomial densities of nonzero total mass: triangulate, restrict
the density to barycentric monomials, and use (4). Each simplex contributes
a polynomial-chain functional; their sum has the original nonzero mass.
For a full-dimensional simplex an affine form is constant on it only if
it is constant on the affine hull, so the point-evaluation exception does
not arise in only some of the triangulation pieces.

\section{Factorial and Gaussian pencils, including coefficient rings}
Let $Y_i$ be independent $\operatorname{Gamma}(\alpha_i,1)$ variables with
positive integer shapes. Put $S=\sum Y_i$ and $L=\sum c_iY_i$.
For a homogeneous binary form $H(U,V)$ of degree $D$, set $F(Y)=H(S,L)$.
Radial separation gives
\[
 \mathbb E[F(Y)^m]=(N)_{Dm}\,
       \mathbb E_{\operatorname{Dir}(\alpha)}[H(1,\ell(T))^m].\tag{5}
\]
Every radial factor is nonzero. Homogeneity recovers $F$ from its restriction
to $S=1$, so Theorem~\ref{thm:simplex} proves:
\begin{corollary}\label{cor:pencil}
Every nonzero $H(S,L)$ has infinitely many nonzero Gamma moments, in every
dimension. With all $\alpha_i=1$ this proves FC on this complete homogeneous
factorial pencil class. Substitution
$Y_i=\frac12\sum_{j=1}^{2\alpha_i}X_{ij}^2$ gives the same conclusion for
the corresponding homogeneous Gaussian phases, retaining the unit marker.
\end{corollary}
This is a homogeneous class using the actual radial sum $S$. It does not
assert FC for an arbitrary inhomogeneous polynomial in $S,L$, or for two
arbitrary linear forms in place of $S,L$. Odd-sized Gaussian blocks and
general nonintegral Dirichlet shapes are not covered by the chain proof.

The field result extends to every characteristic-zero field by embedding
the finitely generated coefficient field into $\C$. Let now $B$ be any
commutative $\Q$-algebra, and allow both the coefficients of $H$ and the
coefficients $c_i$ of $L$ to belong to $B$. Let $J$ be the ideal of the
\emph{actual composed polynomial} $F(Y)=H(S,L)$'s coefficients.
Reduction at every prime and Corollary~\ref{cor:pencil} show that
\[
 \mathbb E_B[F^m]=0\text{ eventually}
 \quad\Longleftrightarrow\quad J\text{ consists of nilpotents}
 \quad\Longleftrightarrow\quad F\text{ is nilpotent}.        \tag{6}
\]
There are only finitely many coefficients, so their nil ideal is nilpotent
without any Noetherian assumption. Nilpotence then gives vanishing against
all original Gamma or Gaussian markers.

The exact ideal of scalars $u$ making every actual marked period polynomial
is
\[
 N_F=\bigcup_{k\ge1}\Ann_B(J^k)=\Gamma_J(B).                \tag{7}
\]
Necessity uses the scaled unit period and the injection
$B/\Ann_B(u)\hookrightarrow B$ given by multiplication by $u$; (6) in
that quotient makes $J^ku=0$ for some $k$. Sufficiency is $uF^m=0$ for
$m\ge k$. This holds at all node collisions and retains ordinary torsion.
Using the coefficients of the formal binary form $H$ instead of $F$ would
be wrong: if $B=\Q[\epsilon]/(\epsilon^2)$, $L=(1+\epsilon)S$ and
$H(U,V)=V-U$, then $F=\epsilon S$ is nilpotent although $H$ has unit
coefficients.

\section{Affine exponential periods at algebraic coefficients}
The next assertion concerns a single exponential value, a different
question from eventual power vanishing.
\begin{theorem}\label{thm:exp}
For positive integer Dirichlet shapes and algebraic $c_i$,
$\mathbb E[\exp(\sum c_iT_i)]\ne0$, including all collisions.
\end{theorem}
\begin{proof}
Repeat the nodes according to the shapes. The unnormalized uniform
simplex integral equals $\sum_\lambda\Res_{z=\lambda}e^z/Q(z)$;
this follows from (2), since $e^z$ is its own iterated derivative.
Write each residue as $a_\lambda e^\lambda$, with $a_\lambda\in\Qb$.
Lindemann--Weierstrass would force all $a_\lambda$ to vanish if their sum
were zero. We show this is impossible.

If every residue of $e^z/Q$ vanished, its local primitives would be
meromorphic. Dividing by $e^z$, choose their principal parts $R_\lambda$,
of pole order at most $m_\lambda-1$. Then $R=\sum R_\lambda$ is rational
and $R'+R-1/Q$ has no finite pole. It tends to zero at infinity, so is
zero. But a nonzero rational $R$ with common denominator of degree at most
$N-s$, where $s\ge1$ is the number of distinct nodes, cannot decay faster
than $z^{-(N-s)}$. The equation $R'+R=1/Q$ forces decay exactly $z^{-N}$,
a contradiction. $R=0$ is impossible as well.

Thus at least one coefficient $a_\lambda$ is nonzero, proving the theorem.
The argument requires algebraic nodes only for Lindemann--Weierstrass.
\end{proof}
Individual coefficients can vanish. For nodes $(0,1,1)$ the integral is
$\int_0^1 z e^z dz=1$, despite a nonconstant affine phase. Thus this theorem
asserts nonvanishing, not transcendence of every nonconstant simplex period.
For transcendental nodes it can fail: nodes $(0,2\pi i)$ give zero.

\section{Flat factorial polynomials with arbitrary lower terms}
The radial expansion of EIT VI extends verbatim to fixed integer Gamma
shapes: for $D\ge2$, $\deg R\le D-2$, and
\[
 f=S^D+D S^{D-2}L+R,
\]
one has
\[
 \lim_{m\to\infty}\frac{\mathbb E[f^m]}{(N)_{Dm}}
       =\mathbb E_{\operatorname{Dir}(\alpha)}[e^{\ell(T)}]. \tag{8}
\]
Indeed the coefficient for a multinomial choice with $K=\sum k_j$ and
$J=\sum jk_j$ is
$\frac{m!}{(m-K)!}\frac{\Gamma(Dm-J+N)}{\Gamma(Dm+N)}$.
It is between zero and one: the first $K$ denominator factors dominate
the numerator and all the others are at least one. Its fixed-index limit
is $D^{-J}m^{K-J}$, and the summable dominating series is
$\prod_j\|A_j\|^{k_j}/k_j!$. Thus only the next-to-top layer survives.

For algebraic $c_i$, Theorem~\ref{thm:exp} makes (8) nonzero, so every
sufficiently large moment is nonzero. The coefficients of $R$ may be
arbitrary complex numbers. With arbitrary complex $c_i$, an all-moment
counterexample in this finite rationally defined family would specialize
to an algebraic one by Noetherianity and Nullstellensatz. The top layer
is fixed and nonzero. Hence FC also holds throughout the complex family.
Taking the $k$th power preserves the displayed form, with degree $Dk$
and the same normalized affine next layer, so eventual pure zero is
also impossible.

For clarity, arbitrary complex coefficients in $L$ give the FC conclusion
and infinitely many nonzero moments; they need not give a nonzero limit
in (8). For two variables, $L=2\pi iY_1$ gives a zero exponential simplex
integral. An overall nonzero scalar multiplying $f$ can be removed before
applying the statement. The degree-one case $f=S+c$ has the direct nonzero
limit $e^c$ after the same radial normalization.

\section{Remaining boundary}
The results include every polynomial of one complex affine coordinate on
a triangle or tetrahedron, with all positive integer Dirichlet weights.
They do not prove the assertion for a general polynomial of two independent
angular coordinates on a triangle, so homogeneous FC(3) remains open here.
Unrestricted SFC and general inhomogeneous FC in three or more variables
also remain unresolved. The source survey's EIT VI/Paper V family audit
and the full bibliography review remain separate work.
\end{document}
